\section{Certify, Then Hedge}
\label{sec:hedge}

The registered gate of Section~\ref{sec:verify} answers one question, whether a
hypothesis may act, and answers it late: it waits until $E_{j,t}$ crosses
$1/\alpha_j$ and then acts on the magnitude the model stated. Our pilot showed both
halves of that design failing in the same place (Section~\ref{sec:results}): the
gate refutes almost nothing that is true, but it activates less than half of the
right-family proposals and a median of ten periods after the proposal, and when it
does act, the stated magnitude is the least reliable field the model produces.
Certify-then-hedge keeps the \eprocess{} and changes what it is used for. The
\eprocess{} no longer switches a controller on; it prices the hypothesis, and the
controller hedges between the baseline and the hypothesis in proportion to that
price.

\subsection{Hypotheses and the language prior}

At each trigger firing $j$ (a stopping time $\tau_j$, at most two per episode) we
register a fixed, action-aligned hypothesis set
$\mathcal{H} = \{\text{demand up}, \text{demand down}, \text{lead-time shift},
\text{demand pulse}\}$. Families on which the compiler has no profitable action
(\texttt{shipment\_loss}, \texttt{transit\_pause}, \texttt{compound}; their
clairvoyant headroom is zero or negative) carry no hypothesis. Hypothesis
$h \in \mathcal{H}$ of firing $j$ receives prior mass $a_j w_{jh}$ with
\begin{equation}
a_j = \alpha\,2^{-j}, \qquad
w_{jh} = \frac{\lambda}{|\mathcal{H}|} + (1-\lambda)\,\Ind\{h = \hat h_j\},
\label{eq:prior}
\end{equation}
where $\hat h_j$ is the family the language model names in its \shockspec{} (none
if it abstains or names a family outside $\mathcal{H}$) and $\lambda\in[0,1]$ is
a registered hedge. The model therefore sets the prior weights and, for the
hypothesis it names, the onset window of its \eprocess{}; nothing else. With
$\lambda = 1$ the arm never calls the model (it is still timed by the shared
trigger, alerts included), and $\lambda = 0$ trusts the model's family alone. The allocation $a_j$ is the
frozen budget of Theorem~\ref{thm:activation}; it is spent on masses rather than
thresholds.

\subsection{Certification, posterior and sizing}

Each $(j,h)$ carries its own \eprocess{} $E_{jh,t}$ of the form
\eqref{eq:eprocess}, started strictly after $\tau_j$ and built from a canonical
specification of $h$ (with the model's onset window when $h = \hat h_j$). Every
hypothesis keeps receiving evidence after a later firing, so nothing is lost to
supersession. With $A_t = \sum_{j:\tau_j < t}\sum_h a_j w_{jh}$ and
$N_t = \sum_{j,h} a_j w_{jh} E_{jh,t}$, the predictable posterior mass of $(j,h)$
at decision $t$ is
\begin{equation}
\pi_{jh,t} = \frac{a_j w_{jh}\, E_{jh,t-1}}{1 - A_t + N_{t-1}},
\qquad
\Pi_t = \sum_{j,h} \pi_{jh,t}.
\label{eq:posterior}
\end{equation}
Validity of what follows needs nothing from the \emph{size} of a shock, so the
size comes from the data rather than the model's magnitude bin. Within a demand
hypothesis a working likelihood of the visible demand path weighs every
registered multiplier, start and duration, and the post-change level is learned
conjugately from an anchor at the pre-onset mean; within the lead-time hypothesis
the same likelihood of visible receipts and the in-transit total weighs offsets
$\Delta \in \{1,2,3\}$ and starts. Each $(j,h)$ thus contributes a predictive law
$G_{jh,t}$ for demand over the protection interval, and the null contributes
arm~1's own $F_{0,t} = \mathcal{N}\bigl((1+L)\bar d_t,\,(1+L)s_{d,t}^2\bigr)$.

\subsection{The hedge}

The order-up-to requirement is the critical fractile $\beta = p/(p+h)$ of the
posterior-predictive mixture of the requirement $Y - P_t$, where $P_t$ is the
in-transit total,
\begin{multline}
x_t = \inf\Bigl\{x : (1-\Pi_t)\,F_{0,t}(x + P_t) \\
+ \sum_{j,h}\pi_{jh,t}\,G_{jh,t}(x + P_t) \ge \beta\Bigr\},
\label{eq:hedge}
\end{multline}
and the order is $q_t = \min(\max(0,\lceil x_t - I_t\rceil), C_t, \hat C_t)$ with
the published smoother $\hat C_t$ unchanged. When no hypothesis is live,
$\Pi_t = 0$ and $q_t$ is arm~1's order bit for bit.

\subsection{Guarantees}

Write $\Prob_0$ for the global no-change null under which every $E_{jh}$ is a
nonnegative supermartingale with unit start (the model-conditional premise of
Theorem~\ref{thm:activation}); this needs the null's baseline to be known.
Hypotheses that can no longer act are dropped from both $N$ and $A$, which
only lowers $\Pi_t$.

\begin{theorem}[Exposure]
\label{thm:exposure}
For every stopping time $\sigma$ and every language-model output, including
abstention, a wrong family or adversarial content,
$\Expect_0[\Pi_\sigma] \le A \le 3\alpha/4$ and, for $c\in(0,1)$,
$\Prob_0(\sup_t \Pi_t \ge c) \le A(1-c)/\bigl(c(1-A)\bigr)$.
\end{theorem}

\emph{Proof sketch.} $N$ is a nonnegative supermartingale with $N_0 \le A$
because the weights are fixed at each $\tau_j$ and sum to at most one, so
optional stopping gives $\Expect_0[N_\sigma] \le A$. Since
$\Pi = N/(1-A+N)$ is concave and increasing in $N$, Jensen's inequality gives the
first bound; Ville's inequality applied to $N/A$ gives the second. \hfill$\square$

\begin{theorem}[Odds regret]
\label{thm:regret}
Let $c_F(x)$ be the one-period newsvendor cost of target $x$ under law $F$, $x^0_t$
arm~1's target and $V_t = c_{G_t}(x^0_t) - \min_x c_{G_t}(x)$ the value of acting
on the hypotheses if they were true. Then
$c_{F_0}(x_t) - c_{F_0}(x^0_t) \le \frac{\Pi_t}{1-\Pi_t}\,V_t$, so under
$\Prob_0$ the expected cumulative surrogate regret is at most
$T A \bar V/(1-A)$.
\end{theorem}

\emph{Proof sketch.} $x_t$ minimizes $(1-\Pi_t)c_{F_0} + \Pi_t c_{G_t}$, so
$(1-\Pi_t)\bigl(c_{F_0}(x_t)-c_{F_0}(x^0_t)\bigr) \le
\Pi_t\bigl(c_{G_t}(x^0_t)-c_{G_t}(x_t)\bigr) \le \Pi_t V_t$; apply
Theorem~\ref{thm:exposure}. \hfill$\square$

The premise is where the guarantee is fragile. Our demand \eprocess{}es, as
registered, estimate the null's mean and spread from the periods before
$\tau_j$. A null simulation of the frozen arm (Section~\ref{sec:limitations})
finds no cell clearly above a bound when that baseline is known, but with the
plug-in estimate and a first firing at period 10, $\Expect_0[\Pi_\sigma]$ reaches
2.0 to 2.7 times $A$; later firings shrink the excess.

Two further properties follow from \eqref{eq:prior} and \eqref{eq:posterior}.
First, with one hypothesis and $\lambda = 0$ the posterior odds are
$a_j E_{j,t}/(1-a_j)$, so thresholding $E_{j,t}$ at $1/\alpha_j$, as the gate
does, is up to one unit of threshold acting when that hypothesis is the more
probable. The correspondence is structural rather than numerical: the gate
certifies the model's full \shockspec{}, magnitude included, while the hedge
certifies a canonical one. Second, the language model can move evidence but not
validity. A hypothesis the model does not name enters with $\lambda$ times its
content-free weight, so a wrong or missing family costs exactly
$\log(1/\lambda)$ nats of evidence against the arm with $\lambda = 1$ ($\log 4$
at $\lambda = 1/4$); the named hypothesis enters with $\lambda + (1-\lambda)|\mathcal{H}|$
times it, a saving of $\log(\lambda + (1-\lambda)|\mathcal{H}|)$ ($\log 3.25$). Neither changes the bounds above. The
theorems are stated for the one-period surrogate that the controller
optimizes; lost sales, the smoother and the zero floor make realized
multi-period cost an empirical matter, which Section~\ref{sec:results} measures.
